\documentclass{article}
\usepackage{graphicx} % Required for inserting images
\usepackage{caption} % Required for \captionof
\usepackage{amsmath}
\usepackage{float}
\usepackage{comment}
\usepackage[T1]{fontenc}
\usepackage[utf8]{inputenc}
\usepackage{lmodern}
\usepackage{microtype}

\title{Hőtani mérés}
\author{Kern Luca, Csongor Vincze}
\date{2026 Április 1}

\begin{document}

\maketitle

\section*{Absztrakt}

Ebben a mérésben hőtani méréseket végeztünk, melyben többek között meghatároztuk a víz fajhőjét elég nagy pontossággal, a hazsnált rozsdamentes acél termoszok hőkapactiitását, illetve más egyebeket. A mérőműszereink közt fontos szerepet játszott a vékony platina ellenállás-hőmérséklet összefüggésén alapuló digitális hőmérő.

\section*{1. Feladat: A hőmérsékleti adatok leolvasása digitális multiméterről}

Ebben a feladatban a multiméteres hőmérsékletleolvasást gyakoroltuk be. A táblázatban szereplő együtthatók és a képlet alapján meghatároztuk, hogy a mért ellenállások mekkora hőmérsékletnek felelnének meg.

Alább találhatóak az adataink, amelyek azt mutatják be, hogy mekkorának mértük a szobahőmérsékletet a két különböző hőmérővel. Mint láthatjuk, ezek az adatok eltérnek, ennek oka az, hogy a ellenállássos hőmérésnél egy eltolást okoz a kapcsolás soros ellenállása hozzáadódik a hőmérő saját ellenállásához. Ezt a következő feladatban tárgyaljuk részletesebben.

zöld termosz:\\
$R_{zold}=110.4 \Omega$\\
$T_{zold}= \text{ 24.5°C}$\\
fekete termosz:\\
$R_{fekete}=111.3 \Omega$\\
$T_{fekete}= \text{ 26.6°C}$\\


\begin{figure}[H]
    \centering
    \includegraphics[width=0.8\textwidth]{1_1_feladat.pdf}
    \caption*{1. feladat: Mért hőmérséklet az idő függvényében}
    \label{}
\end{figure}

Itt a műszer sajnos felvett olyan értékeket is, amik hibásak voltak (overload), ezért azokat kivágtuk a grafikonból.

\section*{2. Feladat: A platinahőmérő kalibrálása}

Ebben a feladatban a hőmérők kalibrálását végeztük el, amire azért volt szükség mert ahogy a hőmérők bekötésre kerültek, a saját ellenállásához hozzáadódott a soros ellenállás is, ami egyy érték beli eltoldódást eredményezett volna a kalibrálás hiányában. Ennek érdekében mindkét termoszba hideg vizet töltöttünk, és egyszerre megmértük a hőmérsékleti egyensúly beállta után a két termoszban a víz hőmérésékletét higanyos hőmérővel, és megmértük az ellenállás-hőmérővel is. A mért hőmérékletek különbségéből könnyen megállappítottuk a soros ellenállásokat.

\begin{table}[H]
    \centering
    \begin{tabular}{|l|c|c|c|}
        \hline
        Termosz & Hőmérővel mért hőmérséklet [°C] & Mért ellenállás [$\Omega$] & Soros ellenállás [$\Omega$] \\ \hline
        Zöld & 22,1 & 110,7 & 1,3 \\ \hline
        Fekete & 22,0 & 111,4 & 2,9 \\ \hline
        Piros & 22,4 & 110,29 & 1,7 \\ \hline
    \end{tabular}
    \caption{Termoszok hőmérsékleti korrekciója}
    \label{tab:hotan_2}
\end{table}

\begin{comment}
Termosz színkódja: zöld (3)\\
Egyensúlyi hőmérséklet: 22,1 °C\\
Mért ellenállás: 110,7 $\Omega$\\
Soros ellenállás: 1,3 $\Omega$\\

Termosz színkódja: fekete\\
Egyensúlyi hőmérséklet: 22,0 °C\\
Mért ellenállás: 111,4 $\Omega$\\
Soros ellenállás: 2,9 $\Omega$\\

Termosz színkódja: piros\\
Egyensúlyi hőmérséklet: 22,4 °C\\
Mért ellenállás: 110,29 $\Omega$\\
Soros ellenállás: 1,7 $\Omega$\\
\end{comment}

Az eredményeinkből látszik, hogy az egyes termoszokban lévő hőmérők valóban kicsit eltérően mérnek. Emelett az is megfigyelhető, hogy a a soros ellenállások csak kis mértékű eltérést okoznak.
Hibaszámítás: .........                  %kell még!

\section*{3. Feladat: A termosz hőveszteségének mérése}

Ebben a feladatban azt mértük meg, hogy ha egy termoszt megtöltünk forró vízzel akkor mekkora a hőveszteség. Ennek a hőveszteségnek két fő oka lehet. Az egyik az, hogy maga a termosz is átmelegszik (ezt a következő feladatban tárgyaljuk részletesebben), illetve hogy a környezetnek is lead a rendszert valamennyi hőt. Mivel a víz és a termosz közti termikus egyensúly beállta után kezdtünk el mérni, így ezzel a méréssel azt tudtukmegállapítani, hogy a rendszert mennyi hőt ad át a környezetének. 5 percig gyűjtöttünk hőmérséklet adatokat, majd ábrázoltuk ezeket az adatokat az eltelt idő függvényében. Erre a grafikonra egyenest illesztettünk, és meghatároztuk a meredekségét, amiből meghatároztuk,hogy egy perc alatt nagyjából mekkora a hőveszteség. Ez $\Delta T=1$ °C-nak adódott.

\begin{figure}[H]
    \centering
    \includegraphics[width=0.8\textwidth]{3_1_graf.pdf}
    \caption*{3. feladat: A termosz hővesztesége az idő függvényében}
    \label{fig:3_1}
\end{figure}

A grafikonról leolvasható, hogy kicsi a hőveszteség, a meredekség szinte 0, ami megfelel a várakozásainknak.
Hőmérsékleti korrekció: .........        %kell még!

\section*{4. Feladat: A termosz hőkapacitásának becslése}

%hibabecslés: nem a tömegmérés relatív hibája, hanem a 2-vel való osztásnak. Annak van egy 10-15-20%-os hibája. Ezt majd át kell gondolni!

Ebben a feladatban a termoszok hőkapacitását próbáltuk megállapítani. Ugyanis mint ahogy azt az előző feladatban is említettük, az hogy a termosz mennyi hőt vesz föl, egy olyan dolog, amely akár nagyban is befolyásolhatja, hogy hogyan alakulnak a hő cserék. Azt próbáltuk megállapítani, hogy mégis mennyire. Megmértük mindkét termosz tömegét, és felltételeztük, hogy a belső ésa külső faluk ugyanolyan vastag,és köztük nincsen számottevő hőcsere, így a termikus folyamatokban  résztvevő tömegeket a teljes tömegek felével közelítettük. Ez után a hőkapacitást úgy számoltuk, hogy az irodalmi adatot használtuk föl a rozsdamentes acél fajhőjére. A mért, számolt és felhasznált adataink alább láthatóak.

$m_{\text{fekete}}=444$ g\\
$m_{\text{zöld}}=430$ g\\
$m_{\text{piros}}=380$ g\\
$c_{\text{rozsdamentes acél}}=500 \text{ J/(kg}^{\circ}\text{C)}$\\
Becsült hőkapacitás:\\
$C_{\text{fekete}}=111$ J/°C\\
$C_{\text{zöld}}=107,5$ J/°C\\
$C_{\text{piros}}=95$ J/°C\\


\section*{5. Feladat: A víz fajhőjének mérése elektromos fűtéssel}

%fontos: előtte és utána is mérni 1-2 percig!!!
Ebben feladatban a víz fajhőjét mértük ki. Ezt a következőképpen tettük: a zöld termoszunkat megtöltöttük hideg vízzel, amit a későbbiekben egy elektromos fűtőtesttel melegítettünk. A mérés kezdetén megmértük a termoszban lévő fűtetlen víz hőméréskletét, tömegét. Ahhoz hogy a fűtőtestet megfelelő teljesítményen tudjuk üzemeltetni (50 W), megmértük az ellenállását. Majd 10 percig folyamatosan úgy melegítettük a termoszban lévő vizet, hogy kevertük az egyenletes elegyedés érdekében.
A mért hőmérséklet adatokat ábrázoltuk egy grafikonon az idő függvényében, amelyre egyenest illesztettünk, és meghatároztuk a meredekségét. A mért adataink alább találhatóak. Ezekből megállapítottuk a víz fajhőjét, amihez a termosz becsült hőkapacitását használtuk fel.

$V_{\text{víz}}=$0,77 L\\
$m_{\text{víz}}=$774 g\\
$R_{\text{fűtőtest}}=1,23 \Omega$\\
$U_{\text{fűtőtest}}=$7,95 V\\
$I_{\text{fűtőtest}}=6,46$ A\\
$T_{kezdeti}=26,4$ °C\\
$T_{veg}=36,4$ °C\\
termosz hőkapacitása az előző feladatból\\
$c_{\text{víz}}=4092 \text{ J/(kg}^{\circ}\text{C)}$\\

Ezek alapján a fajhő $c_{\text{víz}}=4532,52 \text{ J/(kg}^{\circ}\text{C)}$
-> Mi ez az elentmondás?????

Hibaszámítás: ..........



\begin{figure}[H]
    \centering
    \includegraphics[width=0.8\textwidth]{4_1_graf.pdf}
    \caption*{5. feladat: Fűtött víz hőmérsékletfüggése}
    \label{fig:5_1}
\end{figure}

\section*{6. Feladat: Termosz hőkapacitásának mérése keveréssel}

Ebben a feladatban a meleg vizes termosz hőkapacitását mértük ki, olyan módon, hogy a hogy különböző hőmérsékletű folyadékokat öntöttünk össze. Mindkét termoszba beleöntöttük a meleg- illetve hideg vizet, és megmértük a tömegét, majd a termikus egyensúly beállta után a hőméréskeltüket is. Ezután a mágneses keverő használata melett beleöntöttük a hideg vizet a melegvizes termoszba, és közben folyamatosan mértük és feljegyeztük a hőmérséklet adatokat. Ezt követően a kapott adatokból grafikonokat készítettünk, és meghatároztuk a termosz hőkapacitását. A számítások elvégzéséhez feltettük, hogy a rendszer teljesen zártnak tekinhető (nem ad le hőt a környezetének), illetve a víz irodalmi fajhőjét használtuk (szobahőméréskleten).

Ezek az első alkalom eredményei.\\
$m_{hideg}=464$ g\\
$m_{meleg}=467$ g\\
$T_{hideg}=110,4 \Omega$\\
$T_{meleg}=141,0 \Omega$\\
$T_{\text{egyensúlyi}}=125,0 \Omega$\\
$c_{\text{víz}}= 4178 \text{ J/(kg}^{\circ}\text{C)}$ \\  %irodalmi adat
$c_{\text{rozsdamentes acél}}=500 \text{ J/(kg}^{\circ}\text{C)}$\\
Zöld termosz hőkapacitása keveréssel: $C_{termosz}=197.4$\\

Mivel mi piros termoszt kaptuk a második mérési alkalommal, a mérést meg kellett ismételnünk annak ellenére, hogy egy egészen reális (a becsült hőkapacitástól nem teljesen irreális mértékben eltérő) eredményt kaptunk. Ellenben a második alkalommal már közel sem volt annyira jó az eredményünk, mint az első alkalommal. A mérést azért nem ismételtük meg újra, mert a laborvezetőtől azt az utasítást kaptuk, hogy ezt a mérést legfeljebb kétszer végezzük el, és az első alkalom mérésével együtt ez nekünk már a harmadik lett volna.

Ezek a második alkalom eredményei.\\
$m_{hideg}=294$ g\\
$m_{meleg}=406$ g\\
$R_{hideg}=110,29 \Omega$\\
$R_{meleg}=134,45 \Omega$\\
$T_{hideg}=22,3$ °C\\
$T_{meleg}=85,0$ °C\\
$R_{\text{egyensúlyi}}=125,0 \Omega$\\
$T_{\text{egyensúlyi}}=58,4$ °C\\
$c_{\text{víz}}= 4178 \text{ J/(kg}^{\circ}\text{C)}$ \\  %irodalmi adat
$c_{\text{rozsdamentes acél}}=500 \text{ J/(kg}^{\circ}\text{C)}$\\
Piros termosz hőkapacitása keveréssel: $C_{termosz}=23,6$ J/K\\

\begin{figure}[H]
    \centering
    \includegraphics[width=0.8\textwidth]{6_1_graf.pdf}
    \caption*{6. feladat: Zöld termosz hőkapacitásának mérése keveréssel}
    \label{fig:6_1}
\end{figure}

\begin{figure}[H]
    \centering
    \includegraphics[width=0.8\textwidth]{6_2_graf.pdf}
    \caption*{6. feladat: Piros termosz hőkapacitásának mérése keveréssel}
    \label{fig:6_1}
\end{figure}

Hibaszámítás: .........

\section*{7. Feladat Az alumínium fajhőjének mérése}

Ebben a feladatban az volt a célunk, hogy az alumínium fajhőjét határozzuk meg. Ehhez először is megmértük a kapott alumínium henger paramétereit, hogy térfogatának kiszámolásával megbecsülhessük a tömegét. Viszont a továbbiakban a mérleggel mért tömeget használtuk a számolásokban a pontosabb eredményekért. Ezek után mindkét termoszt feltöltöttük annyi vízzel (az egyiket hideggel, a másikat meleggel), hogy ha az alumínium hengert belelógatjuk, akkor telejsen elmerüljön benne, de a víz ne csorduljon ki. A folyadékok tömegét megmértük, majd a hideg vízbe belettük az alumínium hengert, és mindkézt termoszban megmértük a hőmérsékletet, miután beállt a termikus egyensúly. Ezek utána a lehűlt hengert átettük a melegvizes termoszba, és addig mértük és gyűjtöttük a hőmérséklet adatokat, amíg az új termikus egyensúly be nem állt. Ezt ábrázoltuk is egy grafikonon. Ezekből az adatokból pedig kiszámoltuk az alumínium fajhőjét. Megint feltettük azt a számolások során, hogy a környezettel nincsen hőcsere, és megint a víz irodalmi fajhőjét használtuk. A mért és számolt eredményeink alább találhatóak.

$d_{henger}=40.0mm$\\
$h_{henger}=15cm$\\
$V_{henger}= 188.5cm^3$\\
$\rho_{Al}=2,7 \text{ g/cm}^3$\\
$m_{\text{becsült}}=509g$\\
$m_{\text{mért}}=534$\\
hideg termosz adatai:\\
$m_{\text{víz}}=415\text{ g}$\\
$T_{\text{hideg}}=25.45^{\circ}\text{C}$\\
meleg termosz adatai:\\
$m_{\text{víz}}=415\text{ g}$\\
$T_{\text{víz}}=92.2^{\circ}\text{C}$\\
$T_{\text{egyensúlyi}}=75.8^{\circ}\text{C}$\\
$c_{\text{víz}}= 4178 \text{ J/(kg}^{\circ}\text{C)}$ \\  %irodalmi adat
$C_{termosz}=95J/K$\\
$c_{Al}=634 \text{ J/(kg}^{\circ}\text{C)}$\\

\begin{figure}[H]
    \centering
    \includegraphics[width=0.8\textwidth]{7_1_graf.pdf}
    \caption*{7. feladat: Mért hőmérséklet az idő függvényében}
    \label{fig:7_1}
\end{figure}

Az eredményünk ugyan nem lett tökéletes, de nem is volt annyira távol az irodalmi adattól. A legnagyobb hibaforrás valószínűleg az volt, hogy mivel a második alkalommal kapott piros termosz jóval kisebb volt mint a feketék vagy a zöldek, a hőmérő egészen az alumínium henger falához nyomódott, így nem teljesen a víz hőméréskletét mértük, emellett a keverés is nehézkesebben történt meg. Így jóval pontatlanabb volt a hőmérésklet mérés.

Hibaszámítás: .........

\section*{8. Feladat: A sárgaréz fajhőjének mérése}

Ebben a feladatban az előzőhöz nagyon hasonló mérést végeztünk, azzal a különbséggel, hogy most nem alumínium, hanem sárgaréz fajhőjét állapítottuk meg. A réz henger paramétereit ugyanakkoránakmértük, mint az alumínium hengerét. A továbbiakban is megegyezett a mérés menete az előző feladatéval.

$d_{henger}=40mm$\\
$h_{henger}=15cm$\\
$\rho_{henger}=8,73 \text{ g/cm}^3$\\
$m_{\text{becsült}}=1646g$\\
$m_{\text{mért}}=1592$\\
hideg termosz adatai:\\
$T_{\text{hideg}}=26.2^{\circ}\text{C}$\\
meleg termosz adatai:\\
$m_{\text{víz}}=436\text{ g}$\\
$T_{\text{víz}}=82.3^{\circ}\text{C}$\\
$T_{\text{egyensúlyi}}=66.5^{\circ}\text{C}$\\

$c_{\text{víz}}= 4178 \text{ J/(kg}^{\circ}\text{C)}$ \\  %irodalmi adat
$c_{\text{rozsdamentes acél}}=500 \text{ J/(kg}^{\circ}\text{C)}$\\
$C_{termosz}=95 \text{ J/K}$\\
$c_{Cu}=472 \text{ J/(kg}^{\circ}\text{C)}$\\

\begin{figure}[H]
    \centering
    \includegraphics[width=0.8\textwidth]{8_1_graf.pdf}
    \caption*{8. feladat: Mért hőmérséklet az idő függvényében}
    \label{fig:8_1}
\end{figure}

Az eredményeink ezúttal sem lettek tökéletesek, de a rézhengernek már felfüggesztése volt (nem kötél, hanem merev drót), így mikor belehelyeztük a melegvizes termoszba, mégjobban akadályozta a mágneses keverő mozgását. Emellett továbbra is a henger oldalához szorult a hőmérő, így nem pontosan a folyadék hőmérsékletét mértük.

Hibaszámítás: .........

\section*{9. Feladat A jég olvadáshőjének mérése keveréssel}

Ebben a feladatban a jég olvadáshőjét próbáltuk megállapítani olyan módon, hogy a piros termoszunkba meleg vizet töltöttünk, és megvártuk míg ebben a termoszban beáll a termikus egyensúly. Ekkor körülbelül 0$^{\circ}$C-os hőmérsékletű jeget tettünk a meleg vízbe, és elkezdtük gyűjteni az adatokat. Ezeket később grafikon formájában ábrázoltuk. A mérést addig folytattuk, míg a termoszban lévő összes jég el nem olvadt. A mérés végén a termpszban lévő összes víz tömegének megmérésével meghatároztuk a jég tömegét. Ezekből az adatokból pedig kiszámoltuk a jég olvadáshőjét. A számolásokhoz a víz fajhőjére az irodalmi adatot használtuk, és feltettük, hogy a rendszer és a környezet közti hőcsere elhanyagolható. A mért és számolt eredményeink alább láthatóak.

$m_{\text{víz}}=477\text{ g}$ \\
$T_{\text{kezdeti}}=93.2^{\circ}\text{C}$ \\
$T_{\text{egyensúlyi}}=80.8^{\circ}\text{C}$ \\
$m_{\text{teljes}}=510\text{ g}$ \\
$m_{\text{jég}}=33\text{ g}$\\

$c_{\text{víz}}= 4178 \text{ J/(kg}^{\circ}\text{C)}$ \\  %irodalmi adat
$C_{termosz}=95 \text{ J/K}$\\
$L_{f}=450 \text{ kJ/kg}$ \\

\begin{figure}[H]
    \centering
    \includegraphics[width=0.8\textwidth]{9_1_graf.pdf}
    \caption*{9. feladat: Jég olvadáshő}
    \label{fig:9_1}
\end{figure}

Hiba leginkább abból származhatott, hogy a jég hőmérsékletét nem tudtuk pontosan, és nem kizárt, hogy hidegebb volt mint 0$^{\circ}$C, így már az olvadás előtt is hőt vont el a víztől.

Hibaszámítás: .........

\section*{Hivatkozások}

\end{document}